
Fine-tuned vision models often fail to generalize beyond their training benchmarks, yet no model-internal metric exists to detect benchmark-specific encoding before deployment. This work investigates whether benchmark fingerprints---detectable signatures of training dataset origin---exist in model representations and whether they predict cross-dataset performance gaps. Using linear probing on ResNet-50 penultimate layer features, we find that fingerprints are detectable with high accuracy: a logistic regression classifier achieves 99.51\% accuracy (95\% CI: [99.38\%, 99.65\%]) distinguishing models fine-tuned on Flowers102 versus CIFAR-100, with an effect size of Cohen's d = 698.08. However, the Benchmark Fingerprint Score (BFS), defined as classifier confidence for the true benchmark, shows no correlation with generalization gap (Pearson r = 0.022, p = 0.967). All six models achieved BFS > 0.999, creating a ceiling effect that eliminates variance for correlation analysis. These results establish benchmark fingerprints as a measurable phenomenon while demonstrating that fingerprint strength, as measured by BFS, does not predict deployment risk under the tested conditions. The observed paradox---high detectability but no predictive power---suggests that the mechanism linking benchmark concentration to generalization failure may be more complex than simple spurious feature encoding, or that the BFS metric saturates before capturing meaningful variation.

